\chapter*{第二部分学习路线}
\addcontentsline{toc}{chapter}{第二部分学习路线：从力学公式走向守恒思想}
\markboth{第二部分学习路线}{第二部分学习路线}

\begin{learningbox}{从力学公式走向守恒思想}
第二部分是高中力学的主干。前一部分主要回答“物体怎么运动、为什么加速”，这一部分进一步回答“能量怎样转化、动量怎样传递、圆周和天体运动怎样建模”。
\end{learningbox}

\section*{一、先分清三类方法}

\begin{enumerate}[itemsep=0.2em]
  \item \textbf{力和运动方法：} 适合求瞬时加速度、受力关系和约束关系。
  \item \textbf{能量方法：} 适合求速度、位移、最高点、最大压缩量和能量损失。
  \item \textbf{动量方法：} 适合碰撞、爆炸、反冲和相互作用时间很短的问题。
\end{enumerate}

很多综合题难在判断“该用哪一类方法”。读 Day 08 到 Day 14 时，要不断问：题目关注的是力、位移、时间，还是状态变化和守恒条件？

\section*{二、这一部分的认知台阶}

\begin{itemize}[itemsep=0.2em]
  \item Day 08 到 Day 09：从功和能量开始，建立“前后状态比较”的思想。
  \item Day 10：用动量处理相互作用，学会把多个物体看成系统。
  \item Day 11 到 Day 12：把直线运动扩展到圆周和万有引力。
  \item Day 13 到 Day 14：进入振动和波，理解周期性、传播和相位。
\end{itemize}

\section*{三、阅读时特别注意}

看到“光滑、只有重力和弹力做功”，优先想到机械能守恒；看到“碰撞、爆炸、反冲”，先判断系统外力冲量是否可忽略；看到“圆周运动”，先找圆心，再沿半径方向列向心力方程。
